% Terjemahan lengkap Bahasa Indonesia dari chapter1.tex baris 630--762.
% Sumber: Methods of Algebra, Volume 2, commit
% 9a5803ff77dd3257484cb177f851a73770a59dd3 (CC BY 4.0).
% stable-unit-id: o014.aljabr2.chapter1.limits-through-functors

% segment-id: o014.li2.chapter1.limits-through-functors.h001
\section{Meninjau Limit melalui Funktor}\label{sec:limit-functor}
\hypertarget{o014.aljabr2.chapter1.limits-through-functors}{ }

% segment-id: o014.li2.chapter1.limits-through-functors.p001
Berbagai hubungan antara funktor $F: \mathcal{C} \to \mathcal{D}$ dan limit
secara garis besar dapat dibagi ke dalam dua arah. Ambil $\varinjlim$ sebagai
contoh. Diberikan funktor $\alpha: I \to \mathcal{C}$, pertanyaannya adalah:
% segment-id: o014.li2.chapter1.limits-through-functors.l001
\begin{itemize}
	% segment-id: o014.li2.chapter1.limits-through-functors.i001
	\item Andaikan $\varinjlim \alpha$ ada. Apakah citra
	$F\left(\varinjlim \alpha\right)$ memberikan $\varinjlim (F\alpha)$?
	% segment-id: o014.li2.chapter1.limits-through-functors.i002
	\item Andaikan $\varinjlim F\alpha$ ada. Dapatkah limit itu
	``diangkat'' ke $\mathcal{C}$? Dalam arti apa pengangkatan tersebut unik?
\end{itemize}

% segment-id: o014.li2.chapter1.limits-through-functors.p002
Di sini kita terutama membahas kasus $\varinjlim$. Hasil-hasil ini tentu juga
mempunyai versi yang bersesuaian untuk $\varprojlim$; berkat dualitas,
rinciannya tidak perlu diulang.

% segment-id: o014.li2.chapter1.limits-through-functors.p003
Kita akan menata pembahasan $\varinjlim$ dengan bahasa kokonus dan kategori
$(\alpha/\Delta)$.

% segment-id: o014.li2.chapter1.limits-through-functors.c001
\begin{konvensi}\label{con:limit-cone}
	Untuk funktor $\alpha: I \to \mathcal{C}$ yang telah ditetapkan, data
	$(L, (f_i)_{i \in \Obj(I)})$ di $\mathcal{C}$ yang memenuhi syarat-syarat
	berikut disebut \emph{kokonus} dengan alas $\alpha$ dan puncak $L$:
	% segment-id: o014.li2.chapter1.limits-through-functors.l002
	\begin{compactitem}
		% segment-id: o014.li2.chapter1.limits-through-functors.i003
		\item $L$ merupakan objek di $\mathcal{C}$,
		% segment-id: o014.li2.chapter1.limits-through-functors.i004
		\item setiap $f_i : \alpha(i) \to L$ merupakan morfisme di
		$\mathcal{C}$ dan memenuhi syarat kompatibilitas
		$f_j \circ \alpha(i \to j) = f_i$, dengan $[i \to j]$ berkisar pada
		$\mathrm{Mor}(I)$.
	\end{compactitem}
\end{konvensi}

% segment-id: o014.li2.chapter1.limits-through-functors.p004
Dalam terminologi \cite[\S 2.7]{Li1}, kokonus-kokonus ini membentuk suatu
kategori\footnote{Arti tepat $\Delta$ di sini ialah funktor diagonal
$\mathcal{C} \to \mathcal{C}^I$; lihat Contoh
\sourcecrossref{eg:comma-vs-cone}{1.6.3}. Notasi $\Delta$ dipilih dari huruf
awal pinyin \emph{duìjiǎo} (``diagonal''). Namun, mengikuti asas piktografis,
tidak ada salahnya membayangkan $\Delta$ sebagai ``kokonus''; lihat diagram
\eqref{eqn:injlim-diagram}.} $(\alpha / \Delta)$. Secara konkret, setelah alas
$\alpha$ ditetapkan, suatu morfisme dari kokonus $(L, (f_i)_i)$ ke kokonus
$(M, (g_i)_i)$ ialah diagram komutatif
% segment-id: o014.li2.chapter1.limits-through-functors.g001
\begin{equation}\label{eqn:injlim-diagram}\begin{tikzcd}
	& L \arrow[d] & \\
	\alpha(i) \arrow[bend right, rr, "\alpha(i \to j)"'] \arrow[r, "g_i"] \arrow[ru, "f_i"] & M & \alpha(j) \arrow[l, "{g_j}"'] \arrow[lu, "f_j"']
\end{tikzcd}\end{equation}
dengan $i \to j$ berkisar pada $\Mor(I)$; komposisi morfisme didefinisikan
dengan cara yang jelas. Menurut definisi, $\varinjlim \alpha$ beserta keluarga
morfisme kanonik bawaannya
$\iota_i: \alpha(i) \to \varinjlim \alpha$ merupakan objek awal dari
$(\alpha / \Delta)$.
\index{limit}\index[sym1]{lim}

% segment-id: o014.li2.chapter1.limits-through-functors.p005
Funktor $F$ secara langsung menginduksi funktor
$(\alpha/\Delta) \to (F\alpha / \Delta)$, yang memetakan kokonus
$(L, (f_i)_i)$ ke $(FL, (Ff_i)_i)$. Andaikan $\varinjlim \alpha$ ada. Jika
$\varinjlim F\alpha$ juga ada, maka dalam $(F\alpha / \Delta)$ terdapat
morfisme tunggal
% segment-id: o014.li2.chapter1.limits-through-functors.q001
\[ \varinjlim F\alpha \to F(\varinjlim \alpha), \]
dan morfisme ini merupakan isomorfisme jika dan hanya jika
$\left( F(\varinjlim \alpha), (F\iota_i)_i \right)$ memberikan
$\varinjlim F\alpha$.

% segment-id: o014.li2.chapter1.limits-through-functors.p006
Dengan membalik panah-panah, $\varprojlim \alpha$ dapat dicirikan sebagai
objek terminal dari $(\Delta/\alpha)$. Dengan asumsi bahwa kedua limit
projektif yang dibahas ada, dalam $(\Delta/F\alpha)$ terdapat morfisme tunggal
% segment-id: o014.li2.chapter1.limits-through-functors.q002
\[ F\varprojlim \alpha \to \varprojlim F\alpha. \]

% segment-id: o014.li2.chapter1.limits-through-functors.d001
\begin{definisi}\label{def:create-limit}
	\index{limit!mempertahankan, merefleksikan, menciptakan (preserve, reflect, create)}
	Diberikan $F: \mathcal{C} \to \mathcal{D}$ dan funktor
	$\alpha: I \to \mathcal{C}$, tinjau funktor di antara kategori kokonus
	$(\alpha/\Delta) \to (F\alpha/\Delta)$.
	% segment-id: o014.li2.chapter1.limits-through-functors.l003
	\begin{itemize}
		% segment-id: o014.li2.chapter1.limits-through-functors.i005
		\item Funktor $F$ dikatakan \emph{mempertahankan} $\varinjlim$ ini jika
		$(\alpha/\Delta) \to (F\alpha/\Delta)$ memetakan objek awal (jika ada)
		ke objek awal.
		% segment-id: o014.li2.chapter1.limits-through-functors.i006
		\item Funktor $F$ dikatakan \emph{merefleksikan} $\varinjlim$ ini jika,
		apabila suatu objek dari $(\alpha/\Delta)$ dipetakan menjadi objek awal
		dari $(F\alpha/\Delta)$, objek itu sendiri merupakan objek awal.
		% segment-id: o014.li2.chapter1.limits-through-functors.i007
		\item Funktor $F$ dikatakan \emph{menciptakan} $\varinjlim$ ini jika,
		begitu $\varinjlim F\alpha$ ada, $\varinjlim \alpha$ juga ada, dan dalam
		keadaan ini $F$ mempertahankan $\varinjlim$ tersebut dan juga
		merefleksikan $\varinjlim$ tersebut.
	\end{itemize}
	Untuk $\varprojlim$ terdapat konsep-konsep yang bersesuaian. Konsep-konsep
	untuk $\varinjlim$ dan $\varprojlim$ ini dapat saling diperoleh dengan
	mengganti $\mathcal{C}$ dan $\mathcal{D}$ oleh
	$\mathcal{C}^{\opp}$ dan $\mathcal{D}^{\opp}$.
\end{definisi}

% segment-id: o014.li2.chapter1.limits-through-functors.p007
Perhatikan bahwa konsep-konsep ini bergantung pada $\alpha$ yang sedang
dibahas. Sebagai contoh, sangat mungkin $F$ mempertahankan semua
$\varinjlim$ berhingga tetapi tidak mempertahankan $\varinjlim$ secara umum.

% segment-id: o014.li2.chapter1.limits-through-functors.p008
Untuk $\alpha$ tertentu, pernyataan bahwa suatu funktor menciptakan
$\varinjlim$ (atau $\varprojlim$) setara dengan mengatakan bahwa kokonus (atau
konus) di $\mathcal{D}$ yang mendefinisikan limit tersebut dapat diangkat
secara unik ke $\mathcal{C}$, dan hasil pengangkatan itu memberikan limit yang
bersesuaian di $\mathcal{C}$. Keunikan tersebut berasal dari syarat dalam
definisi mengenai merefleksikan $\varinjlim$ (atau $\varprojlim$). Sudut
pandang ini akan tampak jelas dalam Contoh \ref{eg:Grp-Set-create} dan
Proposisi \ref{prop:forget-create} di bawah.

% segment-id: o014.li2.chapter1.limits-through-functors.pr001
\begin{proposisi}\label{prop:ff-reflect}
	Setiap funktor penuh dan setia merefleksikan semua $\varinjlim$ dan
	$\varprojlim$.
\end{proposisi}
% segment-id: o014.li2.chapter1.limits-through-functors.pf001
\begin{bukti}
	Cukup tinjau $\varinjlim$. Jika $F$ penuh dan setia, sifat universal objek
	awal dari $(\alpha/\Delta)$ dapat diverifikasi di $(F\alpha/\Delta)$ setelah
	menerapkan $F$.
\end{bukti}

% segment-id: o014.li2.chapter1.limits-through-functors.p009
Sejalan dengan pembahasan ini, kita perkenalkan sebuah konsep yang berguna.

% segment-id: o014.li2.chapter1.limits-through-functors.d002
\begin{definisi}\label{def:conservative}
	\index{funktor!konservatif (conservative)}
	Funktor $F$ yang memenuhi syarat berikut disebut \emph{konservatif}:
	morfisme $f \in \Mor(\mathcal{C})$ merupakan isomorfisme jika dan hanya jika
	$Ff \in \Mor(\mathcal{D})$ merupakan isomorfisme.
\end{definisi}

% segment-id: o014.li2.chapter1.limits-through-functors.r001
\begin{catatan}\label{rem:conservative-reflect}
	Jika $F$ merupakan funktor konservatif, maka asalkan
	$\varinjlim \alpha$ ada dan $F$ mempertahankan $\varinjlim$ ini, $F$ juga
	harus merefleksikan $\varinjlim$ ini. Alasannya sebagai berikut. Misalkan
	$L \in \Obj((\alpha/\Delta))$ dipetakan menjadi objek awal dari
	$(F\alpha/\Delta)$. Tinjau morfisme tunggal
	$\varinjlim \alpha \to L$ dalam $(\alpha/\Delta)$. Karena $F$
	mempertahankan $\varinjlim$ ini, citra morfisme tersebut di bawah $F$
	harus merupakan isomorfisme. Syarat konservatif kemudian memastikan bahwa
	$\varinjlim \alpha \to L$ juga merupakan isomorfisme.

	Dengan hipotesis ini, syarat merefleksikan $\varinjlim$ dapat dihapus dari
	definisi bahwa $F$ menciptakan $\varinjlim$.
\end{catatan}

% segment-id: o014.li2.chapter1.limits-through-functors.e001
\begin{contoh}\label{eg:Grp-Set-create}
	Kita akan memverifikasi bahwa funktor pelupa $U$ dari kategori grup
	$\cate{Grp}$ ke kategori himpunan $\cate{Set}$ menciptakan semua
	$\varprojlim$ kecil.

	Hal ini merupakan akibat konstruksi konkret $\varprojlim$ dalam
	$\cate{Grp}$ dan $\cate{Set}$; lihat
	\cite[Contoh 2.7.5 dan 2.8.7]{Li1}. Secara lebih terperinci, ambil kategori
	kecil $I$ dan funktor $\beta: I^{\opp} \to \cate{Grp}$. Sebagai himpunan,
	% segment-id: o014.li2.chapter1.limits-through-functors.q003
	\[ \varprojlim U\beta = \left\{\begin{array}{r|l}
		(x_i)_i \in \prod_{i \in \Obj(I)} U\beta(i) & \forall [i \to j] \in \mathrm{Mor}(I), \\
		& \beta(i \to j)(x_j) = x_i
	\end{array}\right\}, \]
	dan proyeksi $p_i: \varprojlim U\beta \to U\beta(i)$ memetakan $(x_j)_j$
	ke $x_i$. Tugas kita sekarang adalah:
	% segment-id: o014.li2.chapter1.limits-through-functors.l004
	\begin{enumerate}
		% segment-id: o014.li2.chapter1.limits-through-functors.i008
		\item melengkapi $\varprojlim U\beta$ dengan struktur grup sedemikian
		hingga setiap $p_i$ merupakan homomorfisme grup, serta menunjukkan bahwa
		struktur grup semacam itu unik;
		% segment-id: o014.li2.chapter1.limits-through-functors.i009
		\item memverifikasi sifat universal $\varprojlim \beta$ bagi grup ini
		beserta keluarga homomorfisme proyeksi $(p_i)_i$.
	\end{enumerate}

	Jelas bahwa $\varprojlim U\beta$ merupakan subgrup dari hasil kali langsung
	$\prod_i \beta(i)$. Kita notasikan subgrup ini sebagai
	$\varprojlim \beta \in \Obj(\cate{Grp})$; inilah satu-satunya struktur grup
	yang membuat setiap $p_i$ menjadi homomorfisme grup. Dengan demikian, tugas
	pertama selesai.

	Selanjutnya, kita verifikasi sifat universal bagi konus
	$\left( \varprojlim \beta, (p_i)_i\right)$. Tinjau sembarang konus
	$(L, (q_i)_i)$ di $\cate{Grp}$, dengan $q_i: L \to \beta(i)$. Sifat
	universal memberikan pemetaan tunggal
	$\varphi: UL \to \varprojlim U\beta$ sedemikian sehingga
	$p_i \varphi = q_i$ berlaku untuk setiap $i$. Sesungguhnya, $\varphi$ secara konkret
	diberikan oleh $y \mapsto (q_i(y))_{i \in \Obj(I)}$, sehingga pemetaan ini
	merupakan homomorfisme grup $L \to \varprojlim \beta$.\footnote{Catatan
	penerjemah: sumber menulis $L \to \varinjlim \beta$; target dikoreksi
	menjadi $L \to \varprojlim \beta$ sesuai limit projektif yang sedang
	dibuktikan (O014-C007).} Verifikasi selesai.
\end{contoh}

% segment-id: o014.li2.chapter1.limits-through-functors.p010
Dengan argumen yang sepenuhnya serupa, dapat ditunjukkan bahwa funktor pelupa
dari kategori ruang Hausdorff kompak $\cate{CHaus}$ ke $\cate{Set}$ juga
menciptakan semua $\varprojlim$ kecil. Verifikasi terkait diserahkan sebagai
latihan dalam bab ini.

% segment-id: o014.li2.chapter1.limits-through-functors.p011
Hasil berikutnya melibatkan kategori funktor berbentuk $\mathcal{C}^J$
(yang mungkin merupakan kategori ``besar'', tetapi ini tidak menjadi masalah).
Kita memerlukan dua langkah persiapan.

% segment-id: o014.li2.chapter1.limits-through-functors.l005
\begin{itemize}
	% segment-id: o014.li2.chapter1.limits-through-functors.i010
	\item Misalkan $J$ merupakan himpunan. Hasil kali $J$ salinan
	$\mathcal{C}$, yakni $\mathcal{C}^J$, didefinisikan sebagai
	$\Obj\left(\mathcal{C}^J\right) = \Obj(\mathcal{C})^J$ dan
	$\mathrm{Mor}\left(\mathcal{C}^J \right) = \mathrm{Mor}(\mathcal{C})^J$.
	Untuk setiap $j \in J$ terdapat funktor proyeksi yang jelas
	$p_j: \mathcal{C}^J \to \mathcal{C}$.

	% Koreksi O014-C008: sumber menggandakan aksara pada frasa "untuk setiap".
	Diberikan kategori $I$ dan funktor $\alpha: I \to \mathcal{C}^J$, jika
	$\varinjlim (p_j \alpha)$ ada untuk setiap $j$, maka
	$\varinjlim \alpha$ juga ada. Konstruksinya dilakukan dengan mengambil limit
	secara ``titik demi titik'' atau ``objek demi objek'':
	$\varinjlim \alpha = \left( \varinjlim (p_j \alpha) \right)_{j \in J}$.

	% segment-id: o014.li2.chapter1.limits-through-functors.i011
	\item Misalkan $J$ merupakan kategori. Kita dapat meninjau kategori funktor
	$\mathcal{C}^J$. Untuk setiap $j \in \Obj(J)$ terdapat funktor evaluasi
	$\mathrm{ev}_j: \mathcal{C}^J \to \mathcal{C}$, yang memetakan objek $F$
	ke $Fj$ dan memetakan morfisme
	$\varphi = (\varphi_{j'})_{j' \in \Obj(J)}$ ke $\varphi_j$. Nanti kita akan
	meninjau funktor berbentuk $\alpha: I \to \mathcal{C}^J$ beserta
	$\varinjlim$-nya. Perhatikan bahwa $\alpha$ dapat dipandang sebagai funktor
	$I \times J \to \mathcal{C}$, dengan nilainya ditulis $\alpha(i, j)$.
\end{itemize}

% segment-id: o014.li2.chapter1.limits-through-functors.p012
Kedua sistem notasi tersebut kompatibel. Jika $J$ merupakan kategori diskret
dan dipandang sebagai himpunan, maka $\mathcal{C}^J$ tidak lain adalah hasil
kali yang dibahas sebelumnya. Untuk $J$ umum, funktor pelupa
$\mathcal{F}: \mathcal{C}^J \to \mathcal{C}^{\Obj(J)}$ memetakan objek $F$ ke
$(Fj)_{j \in \Obj(J)}$ dan memetakan morfisme $\varphi$ ke
$(\varphi_j)_{j \in \Obj(J)}$. Karena itu, $\mathrm{ev}_j = p_j \mathcal{F}$.

% segment-id: o014.li2.chapter1.limits-through-functors.pr002
\begin{proposisi}\label{prop:forget-create}
	Misalkan $J$ dan $\mathcal{C}$ merupakan kategori.
	% segment-id: o014.li2.chapter1.limits-through-functors.l006
	\begin{enumerate}[(i)]
		% segment-id: o014.li2.chapter1.limits-through-functors.i012
		\item Funktor pelupa
		$\mathcal{F}: \mathcal{C}^J \to \mathcal{C}^{\Obj(J)}$ menciptakan
		$\varinjlim$ dan $\varprojlim$.
		% segment-id: o014.li2.chapter1.limits-through-functors.i013
		\item Misalkan $I$ merupakan kategori dan setiap funktor
		$I \to \mathcal{C}$ mempunyai $\varinjlim$ (atau $\varprojlim$).
		Maka semua funktor berbentuk $I \to \mathcal{C}^J$ juga mempunyai
		$\varinjlim$ (atau $\varprojlim$) di $\mathcal{C}^J$.
		% segment-id: o014.li2.chapter1.limits-through-functors.i014
		\item Dengan asumsi pada (ii), untuk setiap $j \in \Obj(J)$, funktor
		evaluasi $\mathrm{ev}_j: \mathcal{C}^J \to \mathcal{C}$ mempertahankan
		$\varinjlim$ (atau $\varprojlim$) untuk diagram berbentuk $I$ tersebut.
		Dengan kata lain, limit di $\mathcal{C}^J$ juga didefinisikan
		secara ``titik demi titik'' atau ``objek demi objek''.
	\end{enumerate}
\end{proposisi}
% segment-id: o014.li2.chapter1.limits-through-functors.pf002
\begin{bukti}
	Cukup tinjau kasus $\varinjlim$. Untuk (i), tetapkan
	$\alpha: I \to \mathcal{C}^J$ dan andaikan $\varinjlim \mathcal{F}\alpha$
	ada serta diberikan oleh kokonus
	% segment-id: o014.li2.chapter1.limits-through-functors.q004
	\[ \left( (X(j))_{j \in \Obj(J)}, \; \left( \iota_i = (\iota_{i, j})_{j \in \Obj(J)} \right)_{i \in \Obj(I)} \right) \]
	dengan $\iota_{i,j}: \alpha(i, j) \to X(j)$. Kita ingin mengangkatnya
	menjadi kokonus di $\mathcal{C}^J$ yang memberikan $\varinjlim \alpha$.

	Pertama, perhatikan bahwa
	$X(j) = \varinjlim \alpha(\cdot, j)$ untuk setiap $j$. Untuk setiap morfisme
	$j \to j'$, sifat funktorial limit (lihat
	\cite[Lema 2.7.4]{Li1}) memberikan morfisme tunggal
	$X(j \to j'): X(j) \to X(j')$ yang membuat diagram berikut komutatif:
	% segment-id: o014.li2.chapter1.limits-through-functors.g002
	\[\begin{tikzcd}
		\alpha(i, j) \arrow[r] \arrow[d, "{\iota_{i, j}}"'] & \alpha(i, j') \arrow[d, "{\iota_{i, j'}}"] \\
		X(j) \arrow[r, "{X(j \to j')}"'] & X(j')
	\end{tikzcd} \quad i \in \Obj(I). \]
	Dengan cara ini, $j \mapsto X(j)$ diperluas menjadi funktor
	$X: J \to \mathcal{C}$, sedangkan
	$\iota_i: \alpha(i, \cdot) \to X$ membentuk morfisme di
	$\mathcal{C}^J$. Tidak sulit memverifikasi bahwa
	$(X, (\iota_i)_i) \in \Obj(\alpha/\Delta)$; selanjutnya kita buktikan
	bahwa objek ini merupakan objek awal. Karena funktor $\mathcal{F}$
	konservatif, syarat merefleksikan $\varinjlim$ tidak perlu diverifikasi lagi.

	Diberikan $(L, (f_i)_i) \in \Obj(\alpha/\Delta)$, sifat universal
	$\varinjlim \mathcal{F}\alpha$ menentukan morfisme tunggal
	$\varphi = (\varphi_j)_j : \mathcal{F}X \to \mathcal{F}L$ sedemikian sehingga
	bagian segitiga pada diagram berikut komutatif untuk setiap $(i, j)$:
	% segment-id: o014.li2.chapter1.limits-through-functors.g003
	\[\begin{tikzcd}
		& X(j) \arrow[d, "\varphi_j"] \arrow[r, "{X(j \to j')}" inner sep=0.8em] & X(j') \arrow[d, "\varphi_{j'}"] \\
		\alpha(i, j) \arrow[ru, "{\iota_{i,j}}"] \arrow[r, "f_{i, j}"' inner sep=0.8em] & L(j) \arrow[r, "{L(j \to j')}"' inner sep=0.8em] & L(j')
	\end{tikzcd}\]
	Jika dapat ditunjukkan bahwa bagian persegi komutatif untuk setiap
	$j \to j'$, maka keluarga komponen $\varphi$ membentuk morfisme
	$X \to L$. Bingkai
	luar sudah diketahui komutatif, sehingga
	% segment-id: o014.li2.chapter1.limits-through-functors.q005
	\[ \varphi_{j'} X(j \to j') \iota_{i, j} = L(j \to j') f_{i, j} = L(j \to j') \varphi_j \iota_{i, j} \]
	untuk setiap $i$. Morfisme dengan domain $X(j)$ ditentukan oleh
	komposisinya dengan semua $\iota_{i,j}$, sehingga bagian persegi tersebut
	komutatif.

	Untuk (ii), tinjau funktor $\alpha: I \to \mathcal{C}^J$. Menurut
	pembahasan sebelumnya, $\mathcal{F}\alpha: I \to \mathcal{C}^{\Obj(J)}$
	mempunyai $\varinjlim$, yang didefinisikan secara objek demi objek. Oleh (i),
	$\mathcal{F}$ menciptakan $\varinjlim \alpha$ di $\mathcal{C}^J$,\footnote{Catatan
	penerjemah: teks sumber menulis $\mathcal{C}$ di sini. Targetnya dikoreksi
	menjadi $\mathcal{C}^J$, yaitu domain funktor pelupa $\mathcal{F}$ dan
	kategori yang memuat funktor $X: J \to \mathcal{C}$ yang baru dibangun;
	lihat koreksi O014-C009.}
	sedangkan pernyataan
	pada (iii) bahwa $\mathrm{ev}_j$ mempertahankan $\varinjlim$ tersebut
	merupakan akibat langsung dari konstruksi ini.
\end{bukti}
